

\subsection{RQ4 (Implication): Do behavioural gaps alias traditional structural metrics?}
\label{sec:rq4}

Line coverage and mutation kill score are the two most common proxies for test suite adequacy~\cite{PAPADAKIS2019275}. Line coverage records whether source lines are executed; mutation kill score additionally requires that assertions detect injected faults. If either metric already captured behavioural gaps, a method with a strong structural score would contain few untested behaviours, and identifying behavioural gaps would add nothing beyond what these metrics already report. 
We test this for each metric in turn by asking whether a method's structural score predicts its behavioural coverage, and whether untested behaviours disappear once that score is high. We report Spearman's $\rho$ (behavioural coverage is not normally distributed)~\cite{spearman1961proof} and the variance explained by OLS $R^2$, and we measure how strongly untested behaviours concentrate in low-scoring methods with Cohen's $h$, comparing the share of UTBs in the top score bucket against the share of methods in that bucket. Table~\ref{tab:rq4} reports the bucket distribution for developer-written suites in the dataset.


\smallskip
\noindent\textbf{Line coverage:} 
Line coverage confers almost no information about behavioural gaps. 
Its correlation with behavioural coverage is negligible ($\rho=0.112, R^2=0.007$).
The distribution of UTBs across coverage buckets nearly matches the distribution of methods: 85.4\% of UTBs fall in the 91–100\% bucket along with 92.3\% of all methods. 
The effect size of the difference is considered \textit{small} (Cohen's $h=0.22$). 
High coverage does not eliminate gaps; among methods with 100\% line coverage, 38.2\% still contain at least one untested behaviour. 
A method can execute every line under test and still leave documented behaviour untested, as demonstrated by the \texttt{emptyArray()} example in Section~\ref{sec:example}.

\smallskip
\noindent\textbf{Mutation kill score:} 
Killing a mutant confers more information because killing a mutant requires assertions to detect behaviours rather than only executing code. 
Kill score correlates more strongly with behavioural coverage ($\rho=0.377$).
%Untested behaviours concentrate in weakly-tested methods: 44.3\% of UTBs fall in the ≤70\% kill score bucket, which holds only 24.3\% of methods 
Untested behaviours are more prevalent in methods with low kill score (44.3\% of UTBs fall in the $\leq70\%$ kill score bucket containing only 24.3\% of methods); the effect size of the relationship is considered \textit{medium} (Cohen's $h=0.41$).
Although a stronger relationship exists than for coverage, kill score explains only 16\% of the variance in behavioural coverage ($R^2=0.162$), leaving most of the difference unaccounted for.

High kill scores also do not eliminate behavioural gaps as 29.6\% of methods with a perfect mutation kill score contain at least one untested behaviour. 
A strong kill score lowers the odds of a behavioural gap without ruling one out. 
For both metrics, a chi-square test of homogeneity confirms that the distribution of UTBs across buckets differs significantly from the distribution of methods ($p<0.001$).

\smallskip
\noindent\textbf{Implication:} 
The same observations above hold for the automatically generated suites. 
Among methods with a perfect structural score, using ASTER 43\% (coverage) / 38\% (kill score) still contain untested behaviours; using EvoSuite, 32\% (coverage) / 34\% (kill score) also contain untested behaviours, so the gap does not close even with generated test suites.

Neither structural metric certifies that a method's documented behaviour is validated. Line coverage is effectively uninformative about behavioural gaps, and mutation kill score is only a weak signal that still misses gaps in nearly a third of methods with a perfect kill score.
Behavioural coverage is therefore a complementary dimension of test adequacy rather than a restatement of either metric. We note that this analysis bounds the relationship from one side only: the 17.5\% figure is a lower bound on confirmed gaps, because \toolname{} is not guaranteed to surface every untested behaviour. 
Whether additional gaps go undetected is left to future evaluation.


\begin{rqbox}{RQ4: Do UTBs alias existing structural metrics?}
Line coverage is effectively uninformative about behavioural gaps, and mutation kill score is only a weak signal: for developer-written tests, 38.2\% of methods with perfect line coverage and 29.6\% with a perfect kill score still contain untested behaviours. This suggests that behavioural coverage is a complementary dimension of test adequacy. %, not a restatement of either metric.
\end{rqbox}


\endinput

% RTH: tried to simplify the discussion above

\textbf{V2:}

Structural metrics such as line coverage and mutation kill score are widely used proxies for test suite adequacy~\cite{PAPADAKIS2019275}. Line coverage measures whether source code lines are executed by tests. Mutation kill score reflects assertion adequacy, as killing a mutant requires tests to both execute the mutated code and detect the behavioural changes through assertions~\cite{10.1145/1062455.1062530}. 
If behavioural gaps are fully captured by these two existing metrics (line execution and assertion adequacy), detecting behavioural gaps would add no additional value.
In this RQ we examine whether behavioural gaps persist independently of these prior metrics.
% For this RQ we aggregate all detected untested behaviours at the method level, since behavioural gaps are identified at the method level. For each method we extract the corresponding mutation score and line coverage from the test report.
%
% \noindent We quantify the relationship between proxy metrics (line coverage and mutation kill score) and behavioural coverage using the following measures for our entire dataset.
\noindent We quantify the relationship between line coverage and mutation kill score and behavioural coverage using three complementary analyses.
% First we compute the Spearman rank correlation~\cite{spearman1961proof} between per-method proxy scores and per-method behavioural coverage percentages. We also report OLS $R^2$ to capture the proportion of variance in behavioural coverage explained by each proxy. 
% Next, we assess predictive power using Nagelkerke pseudo $R^2$~\cite{note1991note} from a univariate logistic regression that predicts binary UTB presence, meaning whether a method has at least one UTB.

\smallskip
\noindent\textbf{Correlation strength and predictive power:} 
We first ask whether a method's line coverage or mutation score correlates with its behavioural coverage, that is the proportion of expected behaviours that are tested by the test suite. A strong correlation would suggest that structural metrics already capture behavioural coverage. 
We use Spearman's rank correlation because behavioural coverage is not normally distributed~\cite{spearman1961proof}. We also report OLS $R^2$ to measure how much of the variance in behavioural coverage each metric explains. 
% We compute the Spearman rank correlation~\cite{spearman1961proof} between per-method proxy scores and per-method behavioural coverage percentages. 
% We also report OLS $R^2$ to capture the proportion of variance in behavioural coverage explained by each proxy. 
To assess how well each metric predicts the \emph{presence} of any UTB in a method, we fit a univariate logistic regression with binary UTB presence as the outcome and compute Nagelkerke pseudo $R^2$~\cite{note1991note}. Low values in these measures would indicate that structural metrics are poor proxies for behavioural coverage.

% We then applied Steiger Z-test~\cite{steiger1980tests} to identify if the correlation between line coverage and behavioural coverage differs significantly from the correlation with mutation score.
% \smallskip
% \noindent\textbf{Comparing the two structural metrics:} 
% We apply the Steiger Z-test~\cite{steiger1980tests} to determine whether the correlation between line coverage and behavioural coverage differs significantly from the correlation between mutation score and behavioural coverage.
Line coverage and mutation score capture different aspects of test quality (execution vs. assertion adequacy). We apply the Steiger Z-test to determine whether the correlation between line coverage and behavioural coverage differs significantly from the correlation between mutation score and behavioural coverage~\cite{steiger1980tests}. A significant difference would reveal which metric, if any, aligns more closely with behavioural completeness.

For developer written test suites, mutation kill score shows a weak monotonic correlation with behavioural coverage (Spearman $\rho=0.377$, OLS $R^2=0.162$). This leaves 83.8\% of the variance in behavioural gaps unexplained by mutation score. 
Line coverage correlates even more weakly ($\rho=0.112$, $R^2=0.007$). 
The Steiger Z-test confirms the difference between these two correlations is statistically significant ($Z=-22.02$, $p<0.001$), indicating that mutation score is a stronger predictor of behavioural coverage compared to line coverage. 
Nagelkerke $R^2$ for predicting UTB presence follows the same pattern: mutation (0.18) and line coverage (0.01).
Both values are low, meaning neither metric reliably predicts the presence of behavioural gaps. This suggests that structural coverage and behavioural coverage capture different aspects of test adequacy.

% Finally, we compute Cohen's $h$~\cite{cohen1988statistical} to measure effect size. This compares how concentrated UTBs are in the high coverage bucket (91\%–100\%) relative to the number of methods falling in that bucket, showing whether high coverage actually corresponds to fewer untested behaviours.
\smallskip
\noindent\textbf{Effect size and significance:} 
% Finally, we compute Cohen's $h$~\cite{cohen1988statistical} to measure the practical effect. 
% This compares the concentration of UTBs in the high coverage bucket (91\%–100\%) against the proportion of methods falling in that bucket, to identify whether high coverage actually corresponds to fewer untested behaviours.
%
Even if correlations are statistically significant, they may be practically irrelevant. We therefore compute Cohen's $h$ to measure the effect size~\cite{cohen1988statistical}. Specifically, we compare the proportion of UTBs that occur within high-coverage methods (91\%–100\%) against the proportion of methods that are high-coverage. A small Cohen's $h$ would mean that high coverage methods have only slightly fewer UTBs than lower coverage methods, confirming that structural metrics do not reliably indicate the absence of behavioural gaps.

Table~\ref{tab:utb-distribution} reports the concentration of UTBs and methods across coverage buckets for developer written test suites, aggregated across all projects. 
%For line coverage, 92.3\% of methods fall in the 91--100\% bucket, but they contain 85.4\% of the identified UTBs. This gap between method concentration and UTB concentration reveals how evenly UTBs are spread across coverage levels. 
Cohen's $h$ quantifies this gap by comparing the proportion of methods in the top bucket against the proportion of UTBs in that bucket. The value is small ($h=0.22$), suggesting high line coverage methods have only a slightly lower UTB density than lower-scoring methods.
%For mutation score, 69.8\% of methods are in the 91--100\% bucket, yet they account for 50.1\% of UTBs, 
With respect to kill score, the value is considered a medium effect ($h=0.41$).
This indicates a clearer separation: methods with higher mutation scores have meaningfully fewer UTBs per method than lower-scoring methods. 
% However, a medium effect means high mutation score does not guarantee the absence of UTBs.

The small/medium effect sizes suggests that UTBs are not strongly concentrated in low‑coverage / kill score methods. They are spread across all levels, including methods with perfect coverage and kill scores. 
UTBs are therefore not an artifact of low coverage or kill scores. 
Among methods with 100\% line coverage, 38.2\% still contain at least one UTB.
%, with 14.6\% of all UTBs appearing in methods below 90\% line coverage. 
For methods with 100\% mutation kill score score, 29.6\% still contain at least one UTB.
%, with nearly half (49.9\%) of all UTBs appearing in methods below 90\% mutation score.
%
A chi‑square test of homogeneity confirms that the distribution of UTBs across buckets differs significantly from the distribution of methods ($p < 0.001$ for both metrics). 
%These results show that behavioural gaps are not just a consequence of unexecuted lines or weak assertions. Instead, behavioural coverage captures a dimension of test adequacy that is largely independent of traditional structural metrics. 
%We further calculated the oracle gap at the method level~\cite{jain2023mind}, finding a positive but weak correlation with UTB count (Spearman $\rho$=0.314, OLS $R^2$=0.076). Yet among the 5,876 methods where the oracle gap is exactly zero, 30.0\% still contain at least one untested behaviour, further confirming that behavioural coverage captures a different dimension of test weakness.

\smallskip
\noindent\textbf{Correlation between structural and behavioural coverage:}


% \smallskip
% \noindent\textbf{UTB distribution across coverage buckets:}\label{par:bucket-behaviour} 

% For developer-written tests, most UTBs appear in the 91--100\% bucket. For mutation testing, most methods have high kill scores, and a large fraction of untested behaviours occur in these methods. In particular, 50.1\% of identified UTBs were found in methods where the mutation score is above 90\%. A similar pattern holds for line coverage, 85.4\% of UTBs were found in methods line coverage above 90\%.
% To further demonstrate that these UTBs are not just from missing line coverage or mutation we also calculated the the percentage of methods having perfect coverage score that still contains at least one UTB (perfect score paradox). We find that 29.8\% of methods with perfect mutation kill score still contains at least one UTB. Similar pattern is observed for line coverage (38.4\% methods have at least one UTB). Both of these evidence shows that behavioural gaps are not just caused by unexecuted lines or weak assertions. This suggests that behavioural coverage captures a different aspect of test adequacy. A chi-square test confirms UTBs are not uniformly distributed across buckets ($p<0.001$). However, Cohen's $h$ compares the proportion of UTBs to the proportion of methods in the 91--100\% bucket to measure how strongly UTBs concentrate at high coverage scores. The effect is only medium for mutation score ($h=0.40$) and small-to-medium for line coverage ($h=0.22$), indicating that high proxy scores do not strongly predict low UTB presence.

% \noindent Overall, neither line coverage nor mutation score provides a strong signal of their absence. This supports the view that behavioural gaps persist even under strong proxy coverage signals.


% \noindent\textbf{Does automatically generated test suites exhibit same pattern?} The same pattern holds for ASTER‑ and EvoSuite‑generated tests (Section~\ref{sec:rq2_5}). Both tools show that behavioural gaps persist under perfect coverage, with 32\%–43\% of perfectly covered methods still having untested behaviours.
The patterns above also hold for the automatically generated test suites (ASTER and Evosuite). 
For ASTER-generated suites, mutation score remains the stronger correlate ($\rho=0.309$ vs. $\rho=0.199$ for line coverage). This also has the perfect-score paradox: 43.3\% of methods with perfect line coverage and 38.1\% with a perfect mutation kill score contain UTBs. 
Cohen's $h$ effect is medium for both metrics ($h=0.27$ for line, $h=0.29$ for mutation). 

EvoSuite generated tests the relationship reverses but the overall conclusion remains similar. Line coverage becomes the stronger correlate ($\rho=0.26$, Nagelkerke $R^2=0.09$) while mutation score drops to near-zero ($\rho=0.049$, Nagelkerke $R^2=0.001$), with a negligible effect size ($h=0.126$). This is expected as EvoSuite generate tests by optimizing for line coverage and can contain weak assertions. Despite this reversal, the perfect paradox still holds: 32.0\% of methods with perfect line coverage and 34.1\% with a perfect mutation kill score still contain UTBs.
% The same pattern holds for automatically generated suites. For ASTER, mutation score remains the stronger correlate ($\rho = 0.309$ vs. $\rho = 0.199$ for line coverage), with medium Cohen's $h$ for both metrics ($h = 0.27$ line, $h = 0.29$ mutation). For EvoSuite, the relationship reverses: line coverage becomes the stronger correlate ($\rho = 0.26$) while mutation score drops to near-zero ($\rho = 0.049$). This is expected given EvoSuite optimizes for line coverage and can produce weak assertions. Despite this reversal, the perfect-score paradox persists across both tools: between 32\%--43\% of methods with perfect line coverage and 34\%--38\% with a perfect mutation kill score still contain UTBs.

\smallskip
\noindent\textbf{Implication:} These results establish that behavioural gaps represent a complementary dimension of test adequacy not captured by either line coverage or mutation kill score. We note that our approach and analysis focuses on precision: the 17.5\% UTB gap in our dataset projects is a lower bound on confirmed gaps, because we cannot claim that \toolname{} identifies all untested behaviours. 
% Quantifying recall would require a complete ground truth of all behaviours for each method, which does not exist. However, for the empirical goal of this study, we believe that establishing that behavioural gaps exist and are prevalent, a precision-focused analysis is sufficient. 
Whether gaps are missed by \toolname remains an open question for future tool evaluation work. 


\begin{rqbox}{RQ4: Do UTBs alias existing structural metrics?}
%\noindent\textbf{Summary:} 
Behavioural gaps have a weak, inconsistent relationship with line coverage and mutation score. 
UTBs persist even when all lines are executed and when mutants are successfully killed.
%, which indicates that both execution-based and assertion-based metrics miss these gaps. This pattern holds for both developer-written and automatically generated test suites. 
Behavioural gaps represent a complementary dimension of test completeness that is not captured by standard structural coverage metrics.
\end{rqbox}
